\subsection{MAKE\_FORMATTER：自定义类型打印}

给自定义类型注册 \texttt{std::formatter<T>}，让它能用于
\texttt{std::format}、\texttt{std::println} 和初始化模板的 \texttt{debug(x)}。
宏定义放在初始化模板；\textbf{注册放在类型定义之后、首次打印之前的全局作用域}。
宏本身需要 C++20 的 \texttt{<format>}，本页的 \texttt{println}、容器打印需要
C++23 的 \texttt{<print>} 与相应标准库支持（\texttt{-std=c++23}）。

\begin{minted}{cpp}
#define MAKE_FORMATTER(T, fmt, ...) \
template <> struct std::formatter<T> { \
    constexpr auto parse(std::format_parse_context& ctx) { \
        return ctx.begin(); \
    } \
    auto format(const T& x, auto& ctx) const { \
        return std::format_to(ctx.out(), fmt, __VA_ARGS__); \
    } \
};
\end{minted}

\textbf{三个参数：}\texttt{T} 是类型，\texttt{fmt} 是格式字符串字面量，
后面的参数是从 \texttt{const T\& x} 取出的成员或表达式。
例如 \texttt{x.a}、\texttt{x.b}、\texttt{x.a + x.b}，顺序对应字符串里的占位符。

\begin{minted}{cpp}
#include <bits/stdc++.h>
using namespace std;
// 上面先放 MAKE_FORMATTER 宏定义
struct Point { int a, b; };
MAKE_FORMATTER(Point, "({}, {})", x.a, x.b)
// 这一句已生成整个 formatter 特化，不用再补分号

int main() {
    Point p{3, 4};
    println("p = {}", p);                // p = (3, 4)
    println("{}", vector<Point>{{3, 4}, {5, 6}});
    // [(3, 4), (5, 6)]
    println(stderr, "{}:{}", "p", p);     // p:(3, 4)，同 debug(p)
}
\end{minted}

\textbf{内部怎么工作：}\texttt{parse} 负责解析冒号后面的格式说明；
这里直接返回 \texttt{ctx.begin()}，表示不消费任何说明，只支持空格式
（\texttt{\{\}} 或 \texttt{\{:\}}）。\texttt{format} 接收对象 \texttt{x}，
通过 \texttt{ctx.out()} 取得输出迭代器，把成员按 \texttt{fmt} 写入目标，
最后返回写完后的迭代器。无须先构造临时字符串。

\textbf{记住这几个边界：}
\texttt{\{:>10\}}、\texttt{\{:.2f\}} 等外层格式不支持；要控制成员格式，
直接把注册字符串写成 \texttt{"(\{:.2f\}, \{:.2f\})"}（成员为浮点数）。
格式串里的字面花括号写成 \texttt{\{\{} 和 \texttt{\}\}}。
同一类型只注册一次，不要再同时使用 \texttt{FMT\_STRUCT}。
带逗号的模板类型先用 \texttt{using} 起别名；仅为自己的类型注册，
不要覆盖标准库已有的 formatter。这一版按原宏保留，至少传一个格式化参数。
